\documentclass[11pt]{article}
\usepackage[margin=1in]{geometry}
\usepackage{amsmath,amssymb}
\usepackage{booktabs}
\usepackage[dvipsnames]{xcolor}
\usepackage{hyperref}
\hypersetup{colorlinks=true,linkcolor=blue!50!black,urlcolor=blue!50!black,citecolor=blue!50!black}

\title{How Novel Is a ``Held-Out'' Reaction?\\
A Within-Training Control for Similarity Audits\\
of the ORDerly Condition-Prediction Benchmark}
\author{David Lindgreen}
\date{October 9, 2026 \\[2pt] \small Version 0.1.0 \quad Preprint, not peer reviewed}

\emergencystretch=3em
\begin{document}
\maketitle

\begin{abstract}
Reaction-condition benchmarks are usually split so that no test reaction shares its exact inputs with training, which is taken to
mean the test set is novel. ORDerly~\cite{wigh2024} follows this practice. We audited its ``reaction string, delete rare'' split
(625{,}697 training and 65{,}445 test reactions) and found that 5.78\% of test products, and 80.84\% of test product scaffolds,
also occur in training, and that 60.5\% of a 1{,}000-row test sample has a training product at Morgan--Tanimoto similarity
$\ge 0.70$. Such numbers are uninterpretable without a reference: how much overlap would any hold-out of this data show? We supply one
by applying the identical measurements to the training set itself, holding out one reaction at a time against rows with a different
input key. The control gives 8.72\% identity overlap, 81.63\% scaffold overlap and 60.2\% at $\ge0.70$. The test split's
similarity to training is therefore what a random hold-out of this redundant corpus produces: the redundancy belongs to the dataset,
not to the released split. The frozen audit and four frequency baselines reproduce exactly. Neural-model results could not be
reproduced because the final checkpoints are not public.
\end{abstract}

\section{Question}
A benchmark's difficulty depends on how different test examples are from training examples. Exact-key separation is necessary but weak:
chemically near-identical products can sit on both sides. The repository asks whether ORDerly's large reported accuracy change
survives an exact reproduction and which data decisions make the benchmark easier. The earlier audit measured test-to-training
overlap but gave no baseline for it. This paper adds that baseline.

\section{Frozen reproduction}
Inputs are the pinned ORDerly benchmark and supplement (checksums committed; the 1.46\,GB archive matched its recorded MD5). The four
top-3 frequency baselines were recomputed from the official reference function and agree with the published figures:
20.24\% and 19.55\% for the reaction-string variants (published about 20\%) and 52.22\% and 51.57\% for the trusted-label
variants (published about 52\%). Re-running the prespecified similarity audit reproduced the committed result with no differences.
\begin{table}[h]
\centering\small
\begin{tabular}{lc}
\toprule
Test-to-training measurement (frozen audit) & Value\\
\midrule
Product identity occurs in training & 5.78\% (3{,}784 of 65{,}444)\\
Product scaffold occurs in training & 80.84\% (51{,}617 of 63{,}852)\\
Sampled max Tanimoto, median / mean & 0.734 / 0.732\\
Sample at Tanimoto $\ge0.70$ & 60.5\% (57.4--63.5)\\
Sample at Tanimoto $\ge0.80$ & 30.0\%\\
Sample at Tanimoto $\ge0.90$ & 11.7\%\\
Sample at Tanimoto $=1.00$ & 8.4\%\\
\bottomrule
\end{tabular}
\caption{Frozen similarity audit; intervals are Wilson 95\%.}
\end{table}

\section{Post-hoc: a within-training control}
For every training reaction we asked the same questions as for a test reaction, using as references only the other training rows whose
exact input key (two reactants and the product string) differs from its own. This mirrors the split's defining property that no test
key appears in training. Canonicalisation, Bemis--Murcko scaffolds and 2{,}048-bit radius-2 Morgan fingerprints are identical to the
frozen audit. The similarity measurement uses a seeded random sample of 1{,}000 training rows. This analysis was not preregistered
and all of it is reported.
\begin{table}[h]
\centering\small
\begin{tabular}{lcc}
\toprule
Measurement & Test vs training & Training vs other training\\
\midrule
Product identity overlap & 5.78\% & 8.72\% (54{,}530 of 625{,}695)\\
Scaffold overlap & 80.84\% & 81.63\%\\
Max Tanimoto, median & 0.734 & 0.729\\
Max Tanimoto, mean & 0.732 & 0.738\\
$\ge0.70$ & 60.5\% & 60.2\% (57.1--63.2)\\
$\ge0.80$ & 30.0\% & 31.8\%\\
$\ge0.90$ & 11.7\% & 14.0\% (12.0--16.3)\\
$=1.00$ & 8.4\% & 11.5\% (9.7--13.6)\\
\bottomrule
\end{tabular}
\caption{Same measurements on the released test split and on the training set held out row by row.}
\end{table}
At the scaffold level and across the similarity distribution up to 0.80 the two columns are indistinguishable. The control is
\emph{higher} than the test split for exact product identity (8.72\% against 5.78\%) and for near-identical similarity
($\ge0.90$ and $=1.00$), so the released split is, if anything, somewhat less redundant with training than an in-sample hold-out. The
mechanism is not determined here; the sampling intervals for the similarity rows are wide enough that only the identity-level
difference is clear.

\section{Discussion}
A high overlap figure does not by itself indict a split. For this corpus, which contains many near-analogue reactions from patents, an
honest random hold-out already has most test scaffolds and about 60\% of products with a close neighbour in training. The finding is
therefore about the \emph{dataset}: accuracy on any split of it partly measures recall of nearby chemistry. That is a statement about
what the benchmark can establish, not a claim that ORDerly leaked information. A split that is substantially more novel than this
control (by scaffold or similarity-based clustering) would be a different and harder benchmark, and the control gives a concrete target
for what ``more novel'' must exceed.

\section{Limitations}
Fingerprint similarity of 1.0 is not molecular identity. The similarity comparison rests on 1{,}000-row samples. Exact-key separation
in the control follows the audit's key definition, not patent family, which is unavailable. Only the ``reaction string, delete
rare'' variant was audited for similarity. The neural-model accuracies could not be reproduced (checkpoints and prediction files are
not public), so nothing here bears on them; no claim is made about prospective performance or wet-lab outcomes.

\section*{Reproducibility and AI assistance}
Code, frozen audits, the control and tests:\\
\url{https://github.com/lindgreendavid/reaction-integrity-lab}\\
Interactive laboratory:\\
\url{https://lindgreendavid.github.io/reaction-integrity-lab/}\\
The analysis code, tests and drafts of this text were produced with substantial assistance from an AI system (Claude, Anthropic);
the author is responsible for the content.

\begin{thebibliography}{9}
\bibitem{wigh2024} D.~S. Wigh, J.~Arrowsmith, A.~Pomberger, K.~C. Felton, A.~A. Lapkin, ORDerly: Data sets and benchmarks for
chemical reaction data, \emph{J. Chem. Inf. Model.} \textbf{64}(9), 3790--3798 (2024), doi:10.1021/acs.jcim.4c00292.
\end{thebibliography}
\end{document}
